% !TeX root = ../教材题库_全部条目_转写工作稿.tex
% Source: XJ-B1-C06-PR-P247-Y0570
% 第6章 / 6.4 用样本估计总体 / 教材 p.247 / PDF 第254页
% OCR未校对：公式与文字均可能有误；未完成数学TeX转写；禁止直接发给学生
% 原题及相邻原解截图（核对用）：../题目原图/XJ-B1-C06-PR-P247-Y0570-01.webp
% 本文件只在转写完成后作为题干使用；OCR 条目尚不是可靠题干。
\begin{bankitem}{p.247 练习 第2题}{第6章 / 6.4 用样本估计总体 / 教材 p.247 / PDF 第254页}{XJ-B1-C06-PR-P247-Y0570}{OCR未校对：公式与文字均可能有误；未完成数学TeX转写；禁止直接发给学生}
\begin{ocrdraft}
[教材p.247 / PDF第254页]
\par
2．一家电池生产厂有A，B两条生产线制造同一型号可充电电池．A，B生
\par
产线的产量比为4：5．现采用分层抽样的方法从某天两条生产线上的成品中随机
\par
抽取样本，并测量产品可充电次数的均值及方差，结果如下：
\par
抽取成品数
\par
样本均值
\par
样本方差
\par
A生产线产品
\par
16
\par
2巧
\par
8
\par
B生产线产品
\par
20
\par
212
\par
13
\par
试根据以上数据计算由36个产品组成的样本的方差，并估计总体方差．
\end{ocrdraft}
\end{bankitem}
